\section{Results}\label{results}

\subsection{Baseline Model}\label{baseline_model}

Table \ref{baseline} indicates the results of the baseline DummyClassifier model:

\begin{table}[htbp]
\caption{Baseline Model Results}
\centering
\resizebox{\columnwidth}{!}{
\begin{tabular}{|c|c|c|c|c|}
    \hline
    \makecell{Train PR-AUC \\ ($\mu \pm \sigma$)} & \makecell{Val PR-AUC \\ ($\mu \pm \sigma$)} & Test PR-AUC & \makecell{Val Rec@P $\geq$ 0.8 \\ ($\mu \pm \sigma$)} & Train time (s)\\
    \hline
    \makecell{0.04868 \\ $\pm$ 0.002375} & \makecell{0.04134 \\ $\pm$ 0.01357} & 0.06050 & \makecell{0.0 \\ $\pm$ 0.0} &  0.9327\\
    \hline
\end{tabular}
}
\label{baseline}
\end{table}

From this table, it can been seen that the baseline model performs on par with the PR-AUC baseline threshold of 4.5\% (equivalent to the ratio of positive cases to the total population). 
The test PR-AUC is slighty higher than the threshold but still inline with what one could expect from a random model that does not use inputs characteristics to predict.

\subsection{Controlled Learning-rate Experiment}\label{controlled_learningrate_section}

The default parameters of the \emph{SGDClassifier} model were used. The only parameters that were changed were the following:

\[
\begin{aligned}
\text{loss} &= \text{log\_loss} \\
\text{learning\_rate} &= \text{constant} \\
\text{class\_weight} &= \text{balanced} \\
\text{random\_state} &= 42.
\end{aligned}
\]

The "class\_weight" parameter is set to 'balanced' to prevent bias towards the majority class. This is especially important when dealing with imbalanced datasets. 
Finally, Table \ref{sgd_learningrate} illustrates the results obtained from sweeping the learning rate under these fixed parameters:

\begin{table}[htbp]
\caption{Learning-rate Sweep Results}
\centering
\resizebox{\columnwidth}{!}{
\begin{tabular}{|c|c|c|c|c|c|}
    \hline
    LR (eta0) & \makecell{Train PR-AUC \\ ($\mu \pm \sigma$)} & \makecell{Val PR-AUC \\ ($\mu \pm \sigma$)} & Test PR-AUC & \makecell{Val Rec@P $\geq$ 0.8 \\ ($\mu \pm \sigma$)} & Train Time (s) \\
    \hline
    1e-4 & \makecell{0.1451 \\ $\pm$ 0.008689} & \makecell{0.09970 \\ $\pm$ 0.03500} & 0.1069 & \makecell{0.0006289 \\ $\pm$ 0.001406} & 1.3874\\
    1e-3 & \makecell{0.1583 \\ $\pm$ 0.01276} & \makecell{0.1070 \\ $\pm$ 0.03744} & 0.1070 & \makecell{0.001258 \\ $\pm$ 0.002813} & 1.2571\\
    1e-2 & \makecell{0.1386 \\ $\pm$ 0.02155} & \makecell{0.09645 \\ $\pm$ 0.03911} & 0.08783 & \makecell{0.001023 \\ $\pm$ 0.001461} & 1.177\\
    \hline
\end{tabular}
}
\label{sgd_learningrate}
\end{table}

\subsection{Grid Search Hyperparameter Tuning}\label{gridsearch_section}

GridSearchCV was then tested to find the most optimal hyperparameters. The following grid space was used:

\[
\begin{aligned}
\text{learning\_rate} &= \{1e-4, 3e-4, 1e-3, 3e-3, 1e-2\} \\
\text{alpha} &= \{1e-6, 1e-5, 1e-4, 1e-3\} \\
\text{penalty} &= \{l_2, l_1\} \\
\text{class\_weight} &= \{None, "balanced"\},
\end{aligned}
\]

together with the default parameters of the \emph{SGDClassifier} model. The only parameters that were changed were the following:

\[
\begin{aligned}
\text{loss} &= \text{log\_loss} \\
\text{learning\_rate} &= \text{constant} \\
\text{random\_state} &= 42.
\end{aligned}
\]

The GridSearchCV function was run with the mentioned grid space on 5 folds. As a result a total of 400 fits were performed with results obtained for each possible configuration.
Table \ref{gridsearch} illustrates the top 5 configurations after running GridSearchCV:

\begin{table}[htbp]
\caption{Top-5 GridSearchCV Results}
\centering
\resizebox{\columnwidth}{!}{
\begin{tabular}{|c|c|c|c|c|c|c|}
    \hline
    Rank & $n_{0}$ & alpha & penalty & class$_{wt}$ & Val PR-AUC ($\mu \pm \sigma$) & Runtime (s)\\
    \hline
    1 & 3e-3 & 1e-4 & l$_{1}$ & balanced & 0.108714 $\pm$ 0.0315008 & 0.3834\\
    2 & 3e-3 & 1e-5 & l$_{1}$ & balanced & 0.108663 $\pm$ 0.0308857 & 0.3844\\
    3 & 3e-3 & 1e-5 & l$_{2}$ & balanced & 0.108534 $\pm$ 0.0309338 & 0.3578\\
    4 & 3e-3 & 1e-6 & l$_{1}$ & balanced & 0.108446 $\pm$ 0.0306997 & 0.4040\\
    5 & 3e-3 & 1e-4 & l$_{2}$ & balanced & 0.108437 $\pm$ 0.0308922 & 0.4040\\
    \hline
\end{tabular}
}
\label{gridsearch}
\end{table}

It was found that the best configuration (Rank 1 in Table \ref{gridsearch}) achieved a Test PR-AUC of 0.1082.

\subsection{HalvingGrid Search Hyperparameter Tuning}\label{halvinggridsearch_section}

HalvingGridSearchCV was then tested to find the most optimal hyperparameters. The same grid used in Section \ref{gridsearch_section} together with the same default parameters were used.
The best hyperparameter configuration was found to be:

\[
\begin{aligned}
\text{eta0} &= 0.0003 \\
\text{alpha} &= 0.0001 \\
\text{penalty} &= \text{l1} \\
\text{class\_weight} &= \text{balanced}.
\end{aligned}
\]

Finally, a summary of the results from the HalvingGridSearchCV can be seen in Figure \ref{halvinggridsearch_table}:

\begin{table}[htbp]
\caption{HalvingGridSearchCV Results}
\centering
\resizebox{\columnwidth}{!}{
\begin{tabular}{|c|c|c|c|}
    \hline
    \#Evals & Runtime (s) & Best Val PR-AUC ($\mu \pm \sigma$) & Test PR-AUC \\
    \hline
    119 & 16.6 & 0.1041 $\pm$ 0.03261 & 0.1093\\
    \hline
\end{tabular}
}
\label{halvinggridsearch_table}
\end{table}

\subsection{Bayesian Optimisation Hyperparameter Tuning}\label{bayesian_section}

Bayesian Optimisation was then run using the default parameters with only the following three changes to ensure consistency:

\[
\begin{aligned}
\text{loss} &= \text{log\_loss} \\
\text{class\_weight} &= \text{balanced} \\
\text{random\_state} &= 42.
\end{aligned}
\]

After running 40 trials using \emph{skopt.gp\_minimize} on the search space seen in Table \ref{bayesian_parameter_searchspace}:

\begin{table}[htbp]
\caption{Parameter Search Space}
\centering
    \begin{tabular}{|c|c|}
    \hline
    Parameter & Value \\
    \hline
    eta0 & 1e-5 - 1e-1 sampled log-uniformly \\
    alpha & 1e-7 - 1e-2 sampled log-uniformly \\
    penalty & {l1, l2} \\
    learning-rate schedule & {constant, optimal, invscaling, adaptive} \\
    \hline
    \end{tabular}
\label{bayesian_parameter_searchspace}
\end{table}

the best hyperparameter configuration was found to be:

\[
\begin{aligned}
\text{eta0} &= 0.009243 \\
\text{alpha} &= 0.0001224 \\
\text{penalty} &= \text{l1} \\
\text{learning-rate schedule} &= \text{adaptive}.
\end{aligned}
\]

In addition, the convergence plot showing the best objective score versus trial number can be seen in Figure \ref{convergence_plot}:

\begin{figure}[h!]
    \centering
    \includegraphics[width=\columnwidth]{images/convergence_plot.png}
    \caption{Convergence Plot}
    \label{convergence_plot}
\end{figure}

Finally, the comparison between the Grid Search hyperparameter tuning method and the Bayesian Optimization hyperparameter tuning method can be seen in Table \ref{bayesian_grid}:

\begin{table}[htbp]
\caption{Bayesian optimisation versus grid search comparison}
\centering
\resizebox{\columnwidth}{!}{
\begin{tabular}{|c|c|c|c|c|}
    \hline
    Method & \#Evals & Runtime (s) & Best Val PR-AUC ($\mu \pm \sigma$) & Test PR-AUC \\
    \hline
    Grid Search & 80 & 27.4 & 0.1087 $\pm$ 0.03150 & 0.1068\\
    Bayesian Opt & 40 & 2.967 & 0.1073 $\pm$ 0.03073 & 0.1102\\
    \hline
\end{tabular}
}
\label{bayesian_grid}
\end{table}
